% 规范第三条：摘要内容（含标题和关键词）**原则上不能超过一页**。
% 故本文件严格控制篇幅：前 3 段给出四问答案，第 4 段给模型与检验，
% 第 5 段给灵敏度/稳健性与创新点，末尾关键词。不再展开推导过程。

\section*{摘\quad 要}
\addcontentsline{toc}{section}{摘要}

本文针对中药材热风烘干过程中的传热传质问题，建立\textbf{一维轴对称径向热湿耦合
数学模型}，采用半控制体有限体积法（FVM）与 L-稳定的隐式时间推进求解，
依次回答四个问题，并以解析基准、收敛阶、守恒收支、退化一致性、灵敏度与稳健性
六类检验验证模型。

\textbf{问题1}（预热平衡，0--1800 s，附录2 常物性）：温度场表面升温快、中心滞后，
1800 s 时中心 $33.58\ ^\circ\mathrm{C}$、表面 $36.79\ ^\circ\mathrm{C}$，
径向温差 $3.21\ ^\circ\mathrm{C}$；水分场干燥完全被限制在表层，
中心含水率 30 min 内\textbf{完全未变}，表面降幅 $40.0\%$。
根因是\textbf{传质特征时间比传热大两个数量级}
（$R_0^2/\alpha\approx2369$ s 对 $R_0^2/D\approx8\times10^4$ s）。
由于附录2 的 $D(C)$ \textbf{不含温度}，问题1 存在\textbf{结构性单向耦合}，
可先解水分、再解温度，无需双向迭代。

\textbf{问题2}（整个烘干过程，附录3 变物性 $+$ 双向耦合）：前 3 h 温度已接近
烘房温度（中心 $49.8494\ ^\circ\mathrm{C}$、表面 $49.9666\ ^\circ\mathrm{C}$），
含水率显著下降但远未达终点（中心 $1.7662$、表面 $1.0081$）。
与问题1 的单向耦合结果相比，双向耦合使 1800 s 中心温度低 $1.386\ ^\circ\mathrm{C}$——
变物性使热扩散率降低 $14\%$，而 $D$ 的 Arrhenius 温度依赖又\textbf{反馈}
压低扩散系数，这正是``问题1 可先水后温、问题2 必须联立''的直接数值体现。

\textbf{问题3}（烘干时间）：按题面``药材\textbf{各处}的水分浓度应低于
$0.15$ kg/kg''，判据必须是\textbf{全域最大值} $\Cmax(t)=\max_r C(r,t)$。
实测 $\Cmax(t)$ 与中心含水率几乎重合，说明中心是最难干燥的位置。
求得 $\tstar = 206959.9521$ s $= \mathbf{57.49}$ \textbf{h} $= 2.3954$ 天。
若误用面积平均判定，会把 2.4 天的工艺说成 1.50 天（\textbf{低估 $37.3\%$}）；
误用表面值则低估 $78.2\%$。

\textbf{问题4}（考虑尺寸变化）：引入材料坐标 $\xi=r/R(t)$，在\textbf{均匀径向仿射
收缩}假设下 $\dot R$ 对流项\textbf{严格为零}（可证，非近似），
且扩散项 $\propto1/R^{2}$ 与表面项 $\propto1/R$ \textbf{不可合并}。
采用附录4 物性，在 $\Delta t=1$ s 的收敛口径下求得
$\tdry = 183926$ s $= \mathbf{51.0906}$ \textbf{h} $= 2.1288$ 天，
终态半径 $1.2000$ cm。
与另一套\textbf{独立实现}（MATLAB、$N=160$、$\Delta t=10$ s 给出 $51.20$ h）
相差 $0.2\%$，构成\textbf{跨实现交叉验证}。
A/B/C 顺序差分显示：换用附录4 物性使干燥\textbf{显著变慢}（固定半径时 $>120$ h
仍未达标），尺寸收缩又\textbf{大幅加速}（$\ge68.9$ h），
净效果为\textbf{加快 $6.40$ h}。

\textbf{模型检验}：Bessel 解析基准最大绝对误差 $<1.2\times10^{-3}\ ^\circ\mathrm{C}$
（相对 $<0.002\%$）；空间收敛阶实测 $+2.20$（理论 $+2$），
时间不低于一阶（\textbf{不声称二阶}）；
内部界面通量成对抵消达机器精度 $1.78\times10^{-16}$；五项退化与一致性检验通过。
$\tstar$ 的误差三分量为事件定位 $0.5$ ms / 时间积分 $1.86$ s /
\textbf{空间离散 $17.4$ s（主导）}，故 $\tstar$ 只报告到 2 位小数。

\textbf{灵敏度与稳健性}：敏感度排序在五档扰动下均为 $D\gg h_m\gg h$，
但 $|S_D|$ 在 $\pm5\%$ 处为 $0.886$、在 $-50\%$ 处升至 $1.794$——
\textbf{灵敏度系数是局部量，报告时必须注明档位}。
数据扰动（删 $10\%$ 数据 / 量化噪声 / 传感器噪声 / 换插值算法）引起的散布为
$0.22\sim67.6$ s，全部经由``长时平台均值''单一渠道传递（$r=-0.9997$）；
而\textbf{平台窗口取法一项即达 $1090$ s，比最大的数据扰动还大 8 倍}——
本文可信度瓶颈在\textbf{建模约定}而非数据质量。

\textbf{主要创新点}：（1）收缩体的材料坐标处理，$\dot R$ 对流项可证为零、
两项 $R$ 幂次不可合并，使问题4 与问题2/3 共用同一套内核并可交叉验证；
（2）提出``\textbf{可忽略性必须绑定判据口径}''的诊断范式——
端面效应在``全域最大值''口径下 $<10^{-5}$、在``体积平均''口径下
$-2.2\%\sim-6.2\%$，\textbf{同一物理效应两种口径差 3 个数量级}。

\vspace{4pt}
\noindent\textbf{关键词}：热风烘干；圆柱体传热传质；有限体积法；水分有效扩散系数；
收缩模型；移动边界；判据口径
